% !TeX program = pdflatex
%==============================================================================
%  Lista Teórica 07 --- Pré-processamento e Pipelines
%  O gabarito sai de "Lista teorica 07 - gabarito.tex".
%==============================================================================
\providecommand{\gabaritoopt}{}
\documentclass[11pt,a4paper]{article}
\usepackage[\gabaritoopt]{../../recursos/latex/estilo-lista}

\begin{document}
\cabecalholista{07}{Pré-processamento e \emph{Pipelines}}

%------------------------------------------------------------------------------
\begin{exercicio}[a pergunta que separa grave de leve]
O critério da Aula~07 para classificar um vazamento não é ``quanto a etapa
aprende dos dados'', e sim uma pergunta só: \textbf{a etapa olha o $Y$?}

Classifique cada situação abaixo como vazamento \textbf{grave}, \textbf{leve} ou
\textbf{nenhum}, e justifique em uma linha aplicando esse critério.
\begin{enumerate}[label=(\alph*),leftmargin=1.1cm]
  \item Padronizar todas as covariáveis usando a média e o desvio-padrão
        calculados no conjunto \emph{completo}, antes de separar treino e teste.
  \item Selecionar, entre $10\,000$ covariáveis, as $50$ mais correlacionadas com
        $Y$ usando todos os dados, e só então rodar a validação cruzada com essas
        $50$.
  \item Ajustar o PCA no conjunto completo, reduzir a $20$ componentes, e só
        então separar treino e teste.
  \item Imputar valores faltantes pela \emph{mediana da coluna}, calculada no
        conjunto completo.
  \item Num banco em que cada paciente aparece em várias linhas (uma por
        consulta), fazer \texttt{KFold} sobre as linhas.
  \item Padronizar dentro de um \texttt{Pipeline}, com \texttt{GridSearchCV} por
        fora.
\end{enumerate}
\end{exercicio}

\begin{solucao}
\begin{center}\small
\renewcommand{\arraystretch}{1.3}
\begin{tabular}{@{}clp{7.4cm}@{}}
\toprule
 & classificação & por quê \\
\midrule
(a) & \textbf{leve}    & O \texttt{StandardScaler} calcula média e desvio das colunas de $\bm X$; nunca vê $Y$. Não há como fabricar sinal a partir de ruído. \\
(b) & \textbf{grave}   & A seleção é feita \emph{pela correlação com $Y$}. Cada dobra de validação já contribuiu para escolher as covariáveis que serão testadas nela. \\
(c) & \textbf{leve}    & O PCA aprende direções de \emph{todas} as colunas ao mesmo tempo, o que parece grave --- mas ele só olha $\bm X$. \\
(d) & \textbf{leve}    & A mediana é calculada na coluna de $\bm X$. Mesmo argumento de (a). \\
(e) & \textbf{grave}   & Não é sobre pré-processamento: é a \emph{mesma unidade} nos dois lados da divisão. O modelo vê consultas do mesmo paciente no treino e na validação. \\
(f) & \textbf{nenhum}  & É exatamente a conduta correta: a padronização é reajustada dentro de cada dobra. \\
\bottomrule
\end{tabular}
\end{center}

Duas observações que a tabela não cabe.

A primeira é sobre (c): ele é o caso em que a intuição mais engana. Ajustar o PCA
no conjunto todo \emph{de fato} usa informação do teste --- as componentes são
outras. Mas ``usar informação'' não é o mesmo que ``inflar o desempenho medido'',
e sem olhar o $Y$ não há como fabricar acerto. A Aula~E2 mede isso.

A segunda é sobre (e): \textbf{nenhum \texttt{Pipeline} protege desse caso}. O
\texttt{Pipeline} garante que cada etapa seja ajustada só no treino da dobra, mas
não tem como saber que duas linhas são o mesmo paciente. A ferramenta é o
\texttt{GroupKFold}, passando o identificador do paciente em \texttt{groups}. Vale
guardar o par: \emph{o \texttt{Pipeline} resolve (a), (c) e (d); só o desenho da
divisão resolve (e)}.
\end{solucao}

%------------------------------------------------------------------------------
\begin{exercicio}[por que um custa 0,40 e o outro custa nada]
A nota traz duas medições. Com $\bm X$ e $y$ gerados \textbf{independentemente}
(isto é, $y$ é ruído puro, e o $R^2$ verdadeiro de qualquer modelo é $0$):
\begin{itemize}[leftmargin=1.4cm]
  \item selecionar as covariáveis mais correlacionadas com $y$ \emph{antes} da
        validação cruzada produz $R^2 = +0{,}40$;
  \item padronizar fora da dobra, num problema com sinal de verdade e escalas
        muito díspares, produz $0{,}8466$ --- contra $0{,}8466$ fazendo certo.
\end{itemize}
\begin{enumerate}[label=(\alph*),leftmargin=1.1cm]
  \item Explique o mecanismo do primeiro número. Se $y$ é ruído puro, de onde sai
        um $R^2$ positivo?
  \item Suponha $d$ covariáveis independentes de $y$, e que você fique com as $m$
        de maior correlação amostral com $y$. O que acontece com o $R^2$ inflado
        quando $d$ cresce, com $n$ e $m$ fixos? Por quê?
  \item Explique por que a padronização não consegue produzir o mesmo efeito, por
        mais que ela use o conjunto todo.
  \item Se o efeito da padronização é desprezível, por que ainda assim se
        recomenda pô-la dentro do \texttt{Pipeline}?
\end{enumerate}
\end{exercicio}

\begin{solucao}
\textbf{(a)} Com $y$ de ruído puro, a correlação \emph{populacional} entre
qualquer covariável e $y$ é zero --- mas a correlação \emph{amostral} não é: ela
flutua em torno de zero com desvio da ordem de $1/\sqrt n$. Ao ficar com as
covariáveis de maior correlação amostral, você está escolhendo, entre milhares de
flutuações, as que por acaso ficaram maiores. Não é um subconjunto qualquer: é o
subconjunto que melhor se ajusta \emph{a este $y$ específico}, inclusive à parte
dele que aparece nas dobras de validação.

Quando a validação cruzada roda depois, cada dobra de validação já participou da
escolha das covariáveis que serão testadas nela. O ruído daquela dobra já está
codificado na lista de variáveis --- e o modelo o reencontra. É o mesmo mecanismo
do Exercício~4 da Lista Teórica~03 (selecionar é otimizar), aqui aplicado a
covariáveis em vez de hiperparâmetros.

\smallskip
\textbf{(b)} O $R^2$ inflado \textbf{aumenta} com $d$. Quanto mais covariáveis
você examina, maior é o máximo das $d$ correlações amostrais --- é o mesmo
argumento do máximo de $M$ variáveis aleatórias do Exercício~4 da Lista~03. Com
$d=10$ você escolhe as $m$ melhores de dez flutuações; com $d=10\,000$, as $m$
melhores de dez mil, que são muito maiores.

Isso é desconfortável, porque a situação em que a seleção de variáveis parece
mais necessária --- $d$ enorme --- é exatamente aquela em que fazê-la errado
custa mais caro.

\smallskip
\textbf{(c)} Porque o \texttt{StandardScaler} \textbf{nunca vê o $y$}. Ele
calcula duas estatísticas por coluna de $\bm X$: a média e o desvio-padrão. Usar
o conjunto completo em vez de só o treino muda essas duas estatísticas por uma
quantidade da ordem de $1/\sqrt n$, e o efeito é o mesmo para todas as
observações, independentemente do valor de $y$ delas.

Dito de outro modo: a etapa não tem nenhum canal por onde a informação sobre a
resposta possa entrar. Ela pode piorar ou melhorar o condicionamento numérico,
mas não pode fazer um modelo parecer melhor do que é. Foi essa a medição:
$0{,}8466$ contra $0{,}8466$.

\smallskip
\textbf{(d)} Três razões, e nenhuma delas é o vazamento:
\begin{enumerate}[label=(\roman*),leftmargin=1.1cm]
  \item \textbf{Custa uma linha.} Se a conduta certa é gratuita, não há motivo
        para debater se a errada é aceitável neste caso.
  \item \textbf{A conduta é a mesma para todas as etapas.} Quem se acostuma a
        padronizar fora do \emph{pipeline} vai, mais cedo ou mais tarde,
        selecionar variáveis fora dele --- e aí custa $0{,}40$. A disciplina vale
        mais que o caso isolado.
  \item \textbf{Reprodutibilidade em produção.} O \texttt{Pipeline} carrega a
        média e o desvio aprendidos no treino, e os aplica automaticamente a
        qualquer observação nova. Sem ele, alguém precisa lembrar de guardar e
        reaplicar esses números --- e é aí que aparecem erros silenciosos.
\end{enumerate}
O que a medição muda não é a conduta; é \textbf{onde gastar vigilância}. Numa
revisão de código, a pergunta cara é ``como as covariáveis foram escolhidas?'' e
``as dobras respeitam as unidades?'', não ``o \texttt{StandardScaler} está no
lugar certo?''.
\end{solucao}

%------------------------------------------------------------------------------
\begin{exercicio}[escrevendo o pipeline]
Um banco tem as seguintes colunas:
\begin{itemize}[leftmargin=1.4cm]
  \item \texttt{idade}, \texttt{renda}, \texttt{saldo} --- numéricas, com valores
        faltantes em \texttt{renda};
  \item \texttt{estado}, \texttt{profissao} --- categóricas, sem faltantes;
  \item \texttt{inadimplente} --- a resposta, binária.
\end{itemize}
A receita desejada é: imputar as numéricas pela mediana, padronizá-las, codificar
as categóricas em variáveis indicadoras, e ajustar uma regressão logística com
penalização $\ell_2$, escolhendo o $C$ por validação cruzada de 5 dobras.
\begin{enumerate}[label=(\alph*),leftmargin=1.1cm]
  \item Escreva o código, usando \texttt{ColumnTransformer}, \texttt{Pipeline} e
        \texttt{GridSearchCV}.
  \item Na sua solução, quantas vezes a mediana de \texttt{renda} é calculada ao
        longo de toda a busca, se a grade de $C$ tem $6$ valores? Justifique.
  \item Por que a ordem imputar $\to$ padronizar importa, e não o contrário?
\end{enumerate}
\end{exercicio}

\begin{solucao}
\textbf{(a)}
\begin{lstlisting}
import numpy as np
import sklearn.model_selection as skm
from sklearn.compose import ColumnTransformer
from sklearn.impute import SimpleImputer
from sklearn.linear_model import LogisticRegression
from sklearn.pipeline import Pipeline
from sklearn.preprocessing import OneHotEncoder, StandardScaler

numericas   = ["idade", "renda", "saldo"]
categoricas = ["estado", "profissao"]

prep = ColumnTransformer([
    ("num", Pipeline([("imp",   SimpleImputer(strategy="median")),
                      ("escala", StandardScaler())]), numericas),
    ("cat", OneHotEncoder(handle_unknown="ignore"), categoricas),
])

tubo = Pipeline([("prep", prep),
                 ("modelo", LogisticRegression(penalty="l2", max_iter=2000))])

busca = skm.GridSearchCV(
    tubo,
    {"modelo__C": np.logspace(-3, 2, 6)},
    cv=skm.StratifiedKFold(5, shuffle=True, random_state=0),
    scoring="roc_auc",
).fit(X, y)
\end{lstlisting}
Dois detalhes que valem nota: \texttt{handle\_unknown="ignore"} evita que o ajuste
quebre quando uma categoria aparece só na validação, e
\texttt{StratifiedKFold} mantém a proporção de inadimplentes em cada dobra --- em
classificação é o padrão do \texttt{GridSearchCV}, mas explicitá-lo deixa o
\texttt{shuffle} sob controle (Exercício~3 da Lista prática~03).

\smallskip
\textbf{(b)} \textbf{30 vezes}: $6$ valores de $C$ $\times$ $5$ dobras. Em cada
combinação, o \texttt{GridSearchCV} chama \texttt{fit} do \emph{pipeline inteiro}
usando só as $4/5$ de treino daquela dobra, e o \texttt{SimpleImputer} recalcula a
mediana ali. (Com \texttt{refit=True}, que é o padrão, há ainda um 31º ajuste no
conjunto completo ao final, para produzir o \texttt{best\_estimator\_}.)

Isso é o oposto de imputar uma vez, no começo, no conjunto todo. Redundante? Sim,
computacionalmente. Mas é o que torna a estimativa honesta --- e note que a
mediana muda de dobra para dobra, o que é justamente a variabilidade que a
validação cruzada precisa capturar.

\smallskip
\textbf{(c)} Porque a padronização usa média e desvio-padrão, e ambos são afetados
pelos valores faltantes. Se você padronizar primeiro, o \texttt{StandardScaler}
encontra \texttt{NaN} na coluna \texttt{renda} e propaga \texttt{NaN} para todas
as estatísticas --- na melhor das hipóteses o código quebra, na pior ele devolve
uma coluna inteira de \texttt{NaN}.

Mesmo que não quebrasse, a ordem inversa seria conceitualmente errada: imputar
depois de padronizar significaria inserir a mediana \emph{da escala padronizada},
o que muda a média da coluna e desfaz a padronização que acabou de ser feita.
Imputar primeiro deixa a coluna completa, e aí a padronização vê os dados que
efetivamente irão para o modelo.
\end{solucao}

%------------------------------------------------------------------------------
\begin{exercicio}[a divisão que nenhum pipeline conserta]
Um laboratório mede a expressão de $20\,000$ genes em $60$ pacientes, três
amostras por paciente ($180$ linhas), e quer prever se o tumor responde ao
tratamento. Um analista monta o \texttt{Pipeline} correto --- imputação e
padronização dentro dele --- e reporta acurácia de $0{,}95$ por validação cruzada
de 10 dobras sobre as $180$ linhas.
\begin{enumerate}[label=(\alph*),leftmargin=1.1cm]
  \item Qual é o problema, e por que o \texttt{Pipeline} não o resolve?
  \item Qual é o número aproximado de pacientes que aparecem \emph{ao mesmo
        tempo} no treino e na validação de uma dobra típica? Justifique.
  \item Como corrigir? Escreva a chamada.
  \item O laboratório também quer selecionar os $200$ genes mais associados à
        resposta. Onde essa etapa tem de ficar?
\end{enumerate}
\end{exercicio}

\begin{solucao}
\textbf{(a)} As três amostras de um mesmo paciente são \textbf{a mesma unidade
experimental}. Ao dividir por linha, quase todo paciente terá amostras no treino e
amostras na validação, e o modelo não precisa aprender nada sobre biologia: basta
reconhecer o perfil de expressão \emph{daquele paciente}, que ele já viu.

O \texttt{Pipeline} não resolve porque ele controla \emph{quando} cada etapa é
ajustada, não \emph{como} as observações são repartidas. Ele garante que a
padronização use só o treino da dobra --- e o treino da dobra já está
contaminado. É a distinção do Exercício~1: o \texttt{Pipeline} resolve os itens
(a), (c) e (d); este é o item (e).

\smallskip
\textbf{(b)} Praticamente \textbf{todos os 60}. Numa dobra de 10 dobras, cerca de
$18$ das $180$ linhas vão para a validação. Cada uma pertence a um paciente que
tem outras duas linhas, e a chance de as outras duas caírem também na mesma
dobra é pequena --- a probabilidade de um paciente ter as três amostras na mesma
dobra de $1/10$ é da ordem de $(1/10)^2 = 1\%$. Logo, quase todos os pacientes
representados na validação também estão no treino, e como são $18$ linhas de até
$18$ pacientes distintos, o vazamento atinge quase toda a dobra.

\smallskip
\textbf{(c)} Dividindo por \textbf{grupo}, com o identificador do paciente:
\begin{lstlisting}
import sklearn.model_selection as skm

cv = skm.GroupKFold(n_splits=10)
notas = skm.cross_val_score(tubo, X, y, cv=cv, groups=id_paciente,
                            scoring="accuracy")
\end{lstlisting}
O \texttt{GroupKFold} garante que todas as linhas de um paciente fiquem do mesmo
lado da divisão. A acurácia vai cair --- e o número menor é o verdadeiro.

\smallskip
\textbf{(d)} \textbf{Dentro do \texttt{Pipeline}}, como um passo
(\texttt{SelectKBest}, por exemplo), de modo que seja reajustada em cada dobra.
Este é o item (b) do Exercício~1, e é o vazamento de $R^2=+0{,}40$ --- agravado
aqui pela proporção: $20\,000$ genes para $60$ pacientes.

Note que o problema deste exercício exige as \textbf{duas} correções ao mesmo
tempo, e elas são independentes. Fazer só o \texttt{GroupKFold} deixa a seleção
de genes vazando; fazer só a seleção dentro do \emph{pipeline} deixa o mesmo
paciente nos dois lados. É por isso que vale ter as duas perguntas na cabeça
separadamente: \emph{a etapa olha o $Y$?} e \emph{a divisão respeita a unidade?}
\end{solucao}


\end{document}
